% Scaffold copied from thesis/planned-toc.md.
% This file intentionally contains headings and local TODO/context comments only.
% Old thesis prose lives in thesis/legacy/ as source material.

\section{Black-Box Data-Science Tools Are Insufficient}\label{sec:black-box-datascience}
%
% Purpose:
% state the negative search result: the attempted black-box and standard
% data-science methods did not produce a second transferable high-`sys` regime.
%
% Writing note:
% negative results can be stated in batches because many rows have the same
% conclusion. The main text can use a table or bullet list.
%
\subsection{Algorithm for random polytopes}\label{subsec:black-box-datascience-algorithm-for-random-polytopes}
%
% State how random polytopes were generated and what search question this tests.
% <!-- OUTLINE GAP: Name the generator, parameters, row counts, retained facets,
%      normalization, and exact negative claim licensed by this sample. -->
%
\subsection{Algorithm for gradient ascent}\label{subsec:black-box-datascience-algorithm-for-gradient-ascent}
%
% State how local ascent was run and what search question this tests.
% <!-- OUTLINE GAP: Name fixed-F/product/continuation variants, seed counts,
%      stopping rules, escape logic, and whether the claim is only local-search
%      negative evidence. -->
%
\subsection{Rows: polytopes}\label{subsec:black-box-datascience-rows-polytopes}
%
% Define the table rows used in the data-science attempts:
% random polytopes, products, ascent endpoints, continuation endpoints, and any
% other retained row families.
% <!-- TOC DECISION: Rows are polytopes. Most rows come from random families and
%      variants of gradient ascent applied to random polytopes; some may come
%      from enumeration instead of randomization. Avoid mixing the non-black-box
%      known HKO2024 `n=1` positive sample into the black-box data-science table,
%      because methods can then memorize that one case without teaching us
%      anything new. -->
% <!-- WRITER-SESSION QUESTION: Finalize row families only when the data-science
%      writer session starts or the branch stabilizes. The set is still dynamic. -->
%
\subsection{Columns: symplectic invariants and metadata}\label{subsec:black-box-datascience-columns-symplectic-invariants-and-metadata}
%
% Define the columns at the level needed to understand the result:
% symplectic quantities, geometric features, orbit data, and metadata.
% <!-- TOC DECISION: Columns are features of the polytopes. They include
%      symplectic invariants, geometric/orbit features, and metadata. Metadata
%      columns are useful caveats because a method may learn data provenance
%      rather than geometry. -->
%
\subsection{Result types}\label{subsec:black-box-datascience-result-types}
%
% Use these result types in the main table or bullet list:
%
% - inapplicable;
% - failed;
% - negative;
% - positive.
%
% Meaning:
%
% - `inapplicable`: the method does not meaningfully apply to the available data
%   or search question.
% - `failed`: the method could not be implemented cheaply enough, scaled enough,
%   or supplied with enough data for a thesis-facing result.
% - `negative`: the method ran and did not produce a useful search rule or new
%   transferable high-`sys` regime.
% - `positive`: the method produced a useful signal or candidate. Use only where
%   the evidence actually supports this.
% <!-- TOC DECISION: If a real positive result appears, report it honestly and
%      escalate to Jorn. It may falsify the current "insufficient" main result
%      and justify spending about another week to follow up toward a better
%      result. Also filter obvious false positives, e.g. a model finding signal
%      because `sys` was regressed against `sys`. -->
%
\subsection{Methods: data-science toolbox}\label{subsec:black-box-datascience-methods-data-science-toolbox}
%
% Batch the attempted, failed, and inapplicable methods. Put detailed figures,
% tables, parameters, and method-specific notes in the data-science appendix.
% <!-- TOC DECISION: Methods come from the data-science toolboxes we actually use
%      and assign to Codex agents. The full set is still dynamic, so the final
%      method list is needed before writing the table, not for the current ToC. -->
% <!-- WRITER-SESSION QUESTION: Finalize the method list during the data-science
%      chapter writer session. For now, keep the ToC generic. -->
%
